% 原文 pp.269–275（印刷 pp.1387–1393）
% 第 19 讲第 2 部分：本文件承接第 1 部分末尾的 §19.2.3，故不写 \section，直接从正文与 \subsubsection 起。
% 说明一：本页范围内没有 Figure/示意图。p.269 的「The N = 2 Hypermultiplet」是排版表格而非插图，
%         故按表格重新排出，未按 §4 裁图（若强行裁图，表格中的数学符号会退化为位图，不符排印规范）。
% 说明二：本页范围内原文带编号的公式只有 (19.2)–(19.4)；其余展示公式原文均无编号，一律用 \[...\]。

% ===== 原文 p.269 =====
% 本页自上一页（p.268）句子中间续起：该句的主语「群」在上一页。

……第 14.4 讲所描述的那个群。把 $\psi_{(2)}$ 与 $\widetilde D=F_A^+-\vec D$ 一并纳入，就给出一个复形，它在形式上计算规范群的等变上同调（equivariant cohomology），取值于伴随值二形式（adjoint-valued two-forms）空间的宇称反转（parity reversal）。由于伴随值二形式的空间是等变可缩的（equivariantly contractible），把它纳入进来并不影响等变上同调。同样，把 $\phi,\psi_{(0)}$ 纳入进来也不改变任何东西，这一点在第 14.4 讲中已经解释过。这就是为什么 $Q$ 上同调中的类总有仅由 $\phi,\psi_{(1)},A$ 构造出来的代表元。

\subsubsection{$\mathcal{N}=2$ 超多重态的扭曲与整体化}
% 原文 §19.2.4：Twisting and Globalizing the N = 2 hypermultiplet

现在我们转而讨论 $\mathcal{N}=2$ 超多重态（hypermultiplet）。下表描述了在扭曲（twisting）之前，这个 $\mathcal{N}=2$ 多重态中各个场的性质。

\begin{center}
表：$\mathcal{N}=2$ 超多重态（原文 The $N=2$ Hypermultiplet，p.269）\\[3pt]
\begin{tabular}{|c|c|c|c|}
\hline
螺旋度（helicity） & $-1/2$ & $0$ & $1/2$\\
\hline
$U(1)_R$ 荷（$U(1)_R$--charge） & $-1$ & $0$ & $1$\\
\hline
$SU(2)_R$ 作用（$SU(2)_R$--action） & 平凡 & $1/2$ & 平凡\\
\hline
场（field） & $\bar\lambda_{\dot\alpha}$ & $M$ & $\lambda_\alpha$\\
\hline
类型（type） & $-$手征旋量 & 玻色子 & $+$手征旋量\\
\hline
\end{tabular}
\end{center}

% [Tutor Note: 读这张表有三处极易出错。其一，$SU(2)_R$ 作用一栏只有 $M$ 取 $1/2$（即 $M$ 是二重态），$\lambda,\bar\lambda$ 都是单态；而拓扑扭曲恰恰是把 $SU(2)_R$ 与某个洛伦兹 $SU(2)$ 对角化，所以只有 $M$ 在扭曲后改变自旋归属（变成 $(1/2,0)$ 的正手征旋量），$\lambda,\bar\lambda$ 原封不动。若误以为三个场都参与扭曲，正、负手征就会整体搞反。其二，本表中 $U(1)_R$ 荷恰为螺旋度的两倍（$-1,0,1$ 对 $-1/2,0,1/2$），这只是该多重态的特殊情形，不能外推到矢量多重态（那里 $U(1)_R$ 荷与螺旋度的关系不同）。其三，$\bar\lambda_{\dot\alpha}$ 上的点号是 Weyl 旋量的手征标记（$\dot\alpha$ 属于 $\bar{\mathbf 2}$），抄录时最易漏掉；一旦漏掉，$S^\pm(X)\otimes\rho$ 这个张量积的旋量类型就再也对不上。]

上表中的玻色子以及两种手征的旋量，都取值于 $X$ 上由表示 $\rho$ 所伴随的四元数丛（quaternionic bundle）。经过扭曲之后，玻色子 $M$ 变成 $(1/2,0)$ 型的旋量，也就是说 $M$ 变成正手征旋量（正自旋丛的一个截面）。费米子 $\lambda$ 与 $\bar\lambda$ 在扭曲下不变，因为 $SU(2)_R$ 平凡地作用于它们。于是，$\lambda$ 仍是正手征旋量，$\bar\lambda$ 仍是负手征旋量。因此，为了写出 Lagrange 量中涉及超多重态的那一部分，我们必须能够理解每一种手征的自旋丛 $S^\pm(X)$ 与四元数表示 $\rho$ 的张量积。（当然，并不要求自旋丛整体存在，只要求这些张量积存在。）

假定我们正处在这种情形，下面描述微分 $Q$ 对这些场的作用。仍把它记为 $\delta$，我们得到

\[
\begin{aligned}
  \delta M&=\lambda,\\
  \delta\lambda&=\phi(M),\\
  \delta\bar\lambda&={\not\partial}_A(M)+\text{fermions}.
\end{aligned}
\]

写 $\phi(M)$ 时，我们的意思是：Lie 代数元素 $\phi$ 按照超多重态所变换的那个规范群表示（无论它是什么表示）作用在超多重态上。

% [Tutor Note: 这三式给出的 $\delta$ 并不幂零。由 $\delta M=\lambda$、$\delta\lambda=\phi(M)$ 立刻得到 $\delta^2M=\phi(M)$，即 $\delta^2$ 是以 $\phi$ 为参数的规范变换（与第 14 讲 p.204 附近的 $\delta^2=[\sigma,\cdot\,]$ 是同一回事）。因此下文所谓「取 $Q$ 的不动点」，应理解为取等变微分的不动点；而 $\delta\bar\lambda$ 只写出领头项 ${\not\partial}_A(M)$、余项笼统记作「$+\text{fermions}$」，是因为完整变换还含费米子的二次项（形如 $\psi_{(1)}\psi_{(1)}$ 的项），这里只保留对确定玻色定点有用的领头项。只保留领头项时，${\not\partial}_A(M)=0$ 只是领头阶的定点方程。另注意 ${\not\partial}_A$ 的下标 $A$：它是斜杠化的协变 Dirac 算符，作用为 $S^+\otimes\rho\to S^-\otimes\rho$（见 p.273）；若漏掉下标、写成普通的 $\not\partial$，则 $\rho$ 非平凡时该式连协变性都不成立。]

\subsection{高能计算的一般形式}
% 原文 §19.3：The general form of the high energy computations

至此，我们已经成功地在任意一个定向 Riemann 四维流形 $X$ 上定义了这个 $\mathcal{N}=2$ 超对称规范理论的扭曲版本——只要该流形上的丛 $S^\pm(X)\otimes\rho$ 有定义即可。

% ===== 原文 p.270 =====

特别地，这包括了任意定向 Riemann 四维流形上的一切纯 $\mathcal{N}=2$ 规范理论。由于理论的紫外行为既与流形上的度量无关，也与任何整体拓扑无关，我们看到：如果我们出发时所取的是一个基础理论（即 $\beta$ 函数为负、从而是渐近自由的理论），那么这个理论在一个完备的 Riemann 四维流形上的实现也将是渐近自由的。因此，在这个理论中计算关联函数是一个提法妥当（well-defined）的问题。

我们关心的是理论中 $Q$ 闭（$Q$-closed）的局域算符的关联函数。有一个局部化（localization）结果，其内容是：要在紫外极限下计算这些关联函数，不必在整个无穷维的场空间上积分，而只需在 $Q$ 作用于场空间的不动点集上积分（见第 14 讲 §4）。换一种方式来想：我们所计算的关联函数与度量无关，因此可以在度量缩为零的极限下计算，或者等价地，由渐近自由，在耦合常数趋于零的极限下计算。在这个极限下，我们做的是关于 Lagrange 量极小点的经典计算——而这些极小点正是 $Q$ 的不动点。

让我们考察 $Q$ 的不动点集的玻色部分。由上面的方程，略去费米子分量并模去规范群的作用，我们看到 $Q$ 的不动点由下式给出：

\[
\begin{aligned}
  d_A\phi&=0,\\
  0&=F_A^+-\vec D,\\
  0&={\not\partial}_A(M).
\end{aligned}
\]

现在，利用运动方程可以把辅助场 $\vec D$ 积掉（integrate out），得到方程

\[
  \vec D=\mu(M),
\]

其中 $\mu$ 是 hyper-Kähler 矩映射（moment map）。于是上面的方程组变成

\begin{equation}
  \tag{19.2}
  \label{eq:19-2}
  d_A\phi=0,
\end{equation}

\begin{equation}
  \tag{19.3}
  \label{eq:19-3}
  F_A^+=\mu(M),
\end{equation}

\begin{equation}
  \tag{19.4}
  \label{eq:19-4}
  0={\not\partial}_A(M).
\end{equation}

在 $A$ 不可约（irreducible，即非可约）的解之外，第一个方程意味着 $\phi=0$，于是我们可以忽略 $\phi$ 以及这个方程。然而在可约解（reducible solutions）处，这个方程是重要的。把 $\phi$ 场包含进来，对于正确处理可约解的贡献极为重要；而可约解当然是 Donaldson 理论中许多重要微妙之处的来源。

% [Tutor Note: 「$A$ 不可约时第一式推出 $\phi=0$」这一步，作者未加任何说明，其实用的是 Schur 引理式的论证：$d_A\phi=0$ 意味着 $\phi$ 与 $A$ 的和乐群（holonomy）在每一点对易；若 $A$ 不可约，其和乐群在 $\mathrm{ad}$ 表示下不可约，则 $\phi$ 只能取中心元素，而对 $SU(2)$（中心为 $\pm1$）再结合 $\phi$ 取 $\mathfrak{su}(2)$ 值（无迹）即得 $\phi=0$。反过来，$\phi\neq0$ 而 $d_A\phi=0$，正是 $A$ 可约（和乐群落在 $\phi$ 的中心化子中）的判据——这就是作者紧接着强调「可约解处这一方程很重要」的确切含义。注意「不可约」这个前提不可省：在平坦联络或非紧情形下，$d_A\phi=0$ 单独还允许 $\phi$ 为非零常数。]

请注意，式 19.3 与 19.4 就是 Seiberg--Witten 方程，只不过这里处于更一般的情形：规范群不必是 $U(1)$。同时注意，如果我们考虑的是纯 $\mathcal{N}=2$ 超对称规范理论（不含任何超多重态），那么 $M=0$，这些方程就化为通常的 Yang--Mills 瞬子方程。

% [Tutor Note: 说 (19.3)(19.4) 是「不要求规范群为 $U(1)$ 的 SW 方程」时，须核对两处对应：标准的 SW 方程写作 $F_A^+=M\overline M$ 与 ${\not\partial}_A M=0$，其中 $M$ 是 $S^+\otimes(\text{某复线丛})$ 的截面；这里 $\rho$ 可以是任意的四元数表示，故 $M$ 是 $S^+\otimes\rho$ 的截面，而右端取 hyper-Kähler 矩映射 $\mu(M)$。当 $\rho$ 平凡（纯 $\mathcal{N}=2$ 理论）时 $M=0$、$\mu=0$，(19.3) 退化为反自对偶方程 $F_A^+=0$，即通常的瞬子方程——这是检验 $F_A^+$ 符号约定（是否含 $\tfrac12(1+*)$ 的投影）最快的办法。若把 $F_A^+$ 与 $F_A^-$ 互换，(19.3) 就变成瞬子而非反瞬子，整套 Donaldson 不变量的符号会整体反过来。]

% ===== 原文 p.271 =====

我们的关联函数将局部化为某些超流形上的积分，而这些超流形底下的几何流形是下述模空间的不交并：

\[
  \mathcal M_E=\left\{
  \begin{aligned}
    &d_A\phi=0\\
    &F_A^+=\mu(M)\\
    &0={\not\partial}_A(M)
  \end{aligned}
  \right\}\Big/(\text{模去规范变换}),
\]

其中并集取遍给定 $E$ 的模空间的所有分支；人们也许会觉得对 $E$ 求和同样自然。

回忆一下：假定 $X$ 紧致，则有一个指标定理可以计算 $\mathcal M_E$ 的虚维数（virtual dimension）。它就是

\[
  \mathrm{virt.\ dim}(\mathcal M_E)
  =k\bigl(4h-c_2(\rho)\bigr)
  -\tfrac12\dim(G)\bigl(\chi(X)+\sigma(X)\bigr)
  -\frac{\sigma(X)}{4}\,\mathrm{quaternion\ dim}(\rho),
\]

其中 $k$ 是丛 $E$ 的瞬子数，$\chi(X)$ 是 $X$ 的 Euler 特征数，$\sigma(X)$ 是它的符号差（signature）。若 $G$ 单连通且单（simple），则

\[
  k=\int_X\frac{\operatorname{Tr}(F\wedge F)}{8\pi^2}.
\]

% [Tutor Note: 上式中的 $4h$ 里 $h$ 是 $G$ 的对偶 Coxeter 数（dual Coxeter number）：$SU(N)$ 上 $h=h^\vee=N$，故 $SU(2)$ 时 $h=2$、$4h=8$，与 p.274 的 $8k$ 相符。核对其它群时务必用对偶 Coxeter 数而不要用 Coxeter 数——对 $B_r,C_r,F_4,G_2$ 二者并不相等（如 $G_2$ 的 Coxeter 数为 $6$、对偶 Coxeter 数为 $4$）。更要紧的是这个系数与 $\beta$ 函数的关系：$\mathcal{N}=2$ 理论的单圈系数正比于 $4h-c_2(\rho)$，所以「$\beta=0$」与「虚维数与 $k$ 无关」是同一件事。作者正是靠这一点把 $\beta$ 的符号翻译成 $k$ 的取值是否有限：$\beta<0$（渐近自由，本丛书取此约定）时虚维数随 $k$ 增大而无界，$\beta>0$ 时随 $k$ 增大而变负，再加上 $k\ll0$ 处的消没定理，才得出「只有有限多个不变量」。若把 $\beta$ 的符号与渐近自由的对应记反，整段论证就会读反。]

% [Tutor Note: 这里的组合是 $\chi(X)+\sigma(X)$，而 p.274 把同一虚维数条件写成 $8k-\tfrac32\bigl(\chi(M)-\sigma(M)\bigr)$；两者在纯 $SU(2)$（$\dim G=3$、$\rho$ 平凡）时理应一致，故其中必有一处印错。用 $\chi=2-2b_1+b_2$、$\sigma=b_2^+-b_2^-$ 可以判定：$\chi+\sigma=2(1-b_1+b_2^+)$ 给出 Atiyah--Hitchin--Singer 的标准结果 $\dim\mathcal M_k=8k-3(1-b_1+b_2^+)$，而 $\chi-\sigma=2(1-b_1+b_2^-)$ 会给出依赖 $b_2^-$ 的结果，与 Donaldson 理论熟知的公式不符。故应以 p.271 的 $+$ 为准，p.274 处疑为笔误。另，p.274 用 $M$ 记四维流形，而全篇（含本页的 $\chi(X)$）用 $X$，且 $M$ 在紧邻的上一段还是超多重态的标量场，读者宜一律按 $X$ 理解。]

请注意，若 $\beta=0$，则 $\mathcal M_E$ 的虚维数与 $E$ 无关，于是（只要这个虚维数非负）$k$ 的一切取值都可以考虑。当 $\beta<0$ 时，随着丛 $E$ 的瞬子数趋于无穷，诸 $\mathcal M_E$ 的虚维数也趋于无穷，因此对于大的 $k$，可以考察看似越来越复杂的关联函数。但当 $\beta>0$ 时，对 $k\gg0$ 虚维数为负。既然有一个消没定理同时表明 $k\ll0$ 时模空间为空，由此可见：$\beta>0$ 时只有有限多个 $k$ 值需要考虑，因而只有有限多个不变量需要计算。这从另一个意义上说明了为什么 $\beta>0$ 的理论是「答案」：$\beta>0$ 理论的拓扑部分只涉及有限多个不变量，而这些不变量有望（对给定的四维流形）原则上在有限多步内算出来；与此相反，$\beta\le0$ 的理论则有着看似无界的复杂性。

正如我们以前说过的，我们感兴趣的是计算形如

\[
  \int_{\Sigma_i}\mathcal O_{P_i}^{(k_i)}
\]

的算符之乘积的期望值，其中 $P_i$ 是 Lie 代数 $\mathfrak g$ 上的不变多项式，$\Sigma_i$ 是 $X$ 中的 $k_i$ 维闭链（$k_i$-cycle）。在纯 $SU(2)$ 规范理论的情形，这些算符以下述方式对应于 Donaldson 理论中通常的算符。由于规范群是 $SU(2)$，不可约多项式构成的空间由 $\operatorname{Tr}\phi^2=u$ 生成。若我们用这个多项式不变量来定义 $\mathcal O^{(0)}$，则这个算符对应于 Donaldson 理论中通常记作 $\mu(\mathrm{pt})$ 的东西。也就是说，该算符对应于模空间上的这样一个上同调类：它是模空间上由「带基点模空间」（based moduli space）给出的 $SO(3)$ 丛的第一 Pontrjagin 类的负四分之一。此外，导出算符（derived operators）

\[
  \int_\Sigma\mathcal O^{(k)}
\]

（其中 $\Sigma$ 是一个 $k$ 维闭链）对应于上同调类 $\mu([\Sigma])\in H^{4-k}(\mathcal M_E)$。

% ===== 原文 p.272 =====

更一般地，我们可以从 Lie 代数上任一个次数为 $d$、在伴随表示下不变的多项式函数 $P$ 出发，造出 $\mathcal O_P^{(0)}$ 以及导出算符 $\mathcal O_P^{(n)}$。一般而言，由拓扑上的考虑（Künneth 公式）我们知道，如何把由一个可分解（decomposable）不变多项式函数导出的算符，分解为若干由不可分解不变函数构造的算符之乘积的和。因此，一般而言只需同由 Lie 代数上不可分解的不变多项式函数导出的算符打交道就够了。在 $\mathfrak{su}(2)$ 的情形，唯一的不可分解不变多项式是 $\operatorname{Tr}\phi^2$。因此，只需计算由这一个零维局域算符导出的算符。只限于这个算符及其「后代」（descendants），在 Donaldson 理论中是标准做法。

一般而言，为了计算关联函数

\[
  \left\langle\prod_{i=1}^{s}\int_{\Sigma_i}\mathcal O_{P}^{(b_i)}\right\rangle,
\]

其中 $P_i$ 的次数为 $d_i$，我们必须对解的超流形的所有分支积分，而这些分支底下的几何流形 $\mathcal M_E$ 满足

\[
  \mathrm{virt.\ dim}(\mathcal M_E)=\sum_{i=1}^{s}(2d_i-b_i).
\]

事实上，当 $G$ 不单连通时，我们还可以做得稍好一些。我们可以对所有这样的主 $G$ 丛 $E$ 求和：它们在 $X$ 的二骨架（two-skeleton）上具有固定的同构类。\footnote{若 $G$ 单连通，则这个改进是空洞的：因为此时所有 $G$ 丛限制到二骨架上都是平凡丛。}这个限制与路径积分的计算「可交换」（即相容），办法是把路径积分用低能有效 Lagrange 量来表达，因为丛在二骨架上的限制不可能通过改变 $X$ 中某个球内的数据而改变。\footnote{当 $X$ 的度量被放大时，该球逼近 $\mathbb{R}^4$，而球内的理论可以借助 $\mathbb{R}^4$ 上的低能有效场论来分析。低能有效场论把球内所有场的效果都求和进来，这些场在球的边界上取给定的值（相对低能的值）。所有这样的场在二骨架上都有相同的拓扑限制。}一个更显然的限制，即固定 $k$ 的值，就不会如此行为良好；因为在时空的每个局部区域中对瞬子数的一切取值求和，正是我们上周引入的低能有效 Lagrange 量之构造中的要素之一。

% [Tutor Note: p.272 的两个脚注（原文编号 2、3）合起来说明了「固定 $E$ 在二骨架上的同构类」为何与低能有效 Lagrange 量相容。拓扑上的要点是：从一个四维流形中挖去一个四维球并不改变它的二骨架（$B^4$ 的二骨架缩为一点），所以在球内改动瞬子数，丛在二骨架上的限制丝毫不变；而固定 $k$ 之所以不行，是因为低能有效理论在构造时已经把每个局部区域内的所有瞬子数求和吸收进了有效耦合。还要记住四维流形上丛的拓扑分类：同构类由「二骨架上的限制」加上「四维胞腔上的瞬子数（即 $c_2$ 的积分）」共同决定——这正是文中在固定 $k$ 之外还须固定二骨架上之类的原因，二者缺一不可。]

当所有瞬子都光滑且不可约时，瞬子解的超流形就是通常瞬子模空间的切丛的宇称反转。在那种特殊类型的超流形上对函数积分，等价于在底下的几何流形上对微分形式积分。维数方程保证我们所积分的是超流形底下那个几何流形 $\mathcal M_E$ 上的最高维形式。

\subsubsection{$\beta=0$ 的一个情形}
% 原文 §19.3.1：A case when β = 0

让我们继续讨论 $\beta=0$ 的情形。为简单起见，取规范群 $G$ 为连通的（紧）半单群。于是，由于诸 $\mathcal M_E$ 都有相同的形式维数，我们的积分是一次性地对瞬子数的一切取值求和。为了精确地写出正确的公式，我们需要下面这个微妙之处。

% ===== 原文 p.273 =====

我们并未在弯曲的四维流形上显式写出 Lagrange 量，只写出了平直四维空间上的。把理论放到弯曲空间上时，我们可以向作用量加入一些涉及度量的项。由于 $\beta=0$，这些理论有一个定义良好的耦合常数 $\tau$。保持拓扑不变性的这样的项有两个，即

\[
  \int_X\bigl(f(\tau)\mathrm{Pf}(R)+g(\tau)L(R)\bigr),
\]

其中 $R$ 是四维流形上度量的 Riemann 曲率张量，$\mathrm{Pf}$ 表示 Pfaffian，$L$ 表示 $L$ 多项式（$L$-polynomial），而 $f,g$ 是 $\tau$ 的全纯函数。（由于 $\bar\partial/\partial\bar\tau$ 作用于 Lagrange 量给出 $\{Q,\cdot\}$，任何 $Q$ 不变算符之乘积的期望值都是 $\tau$ 的全纯函数。）这两项给出拓扑项，分别计算 $\chi(X)$ 与 $\sigma(X)$ 的普适倍数。把这些项加进来，相当于在关联函数中插入一个乘性因子。即使我们出发时所用的经典 Lagrange 量不含这些额外项，它们在重整化过程中也会被加进来，而它们出现的系数将依赖于该过程的具体细节。因此，最终我们所计算的东西具有如下形式：

\[
  q^{a(\tau)\chi(X)+b(\tau)\sigma(X)}
  \sum_k\left\langle\prod_{i=1}^{s}\int_{\Sigma_i^{b_i}}\mathcal O^{(b_i)}\right\rangle_k q^k,
\]

其中 $a,b$ 是由 $f,g$ 导出的函数。例如，若 $f,g$ 是 $\tau$ 的线性函数，则 $a,b$ 是常数。这里，我们已经固定了 $E$ 在二骨架上的同构类，并用 $k$ 表示瞬子数。\footnote{视规范群而定，瞬子数可能不是整数；但它的模 $\mathbb{Z}$ 之值由丛在二骨架上的同构类所决定。}此外，$q=e^{2\pi i\tau}$ 是通常的瞬子计数参数，而因子 $q^k\cdot q^{a(\tau)\chi(X)+b(\tau)\sigma(X)}$ 就是 $\exp(-S)$ 在瞬子数为 $k$ 的那一族经典解的模空间分支上的取值。由于这类理论被认为在 $SL(2,\mathbb{Z})$ 的一个有限指数子群 $\Gamma$ 下不变，这个结果应当在 $\Gamma$ 下像一个模形式（modular form）那样变换，可能在上半平面被 $\Gamma$ 所取的商空间的尖点（cusps）处带有极点。

% [Tutor Note: 括号里的 $\bar\partial/\partial\bar\tau$ 是物理文献的惯用写法，意思是「关于 $\bar\tau$ 的导数」（$\partial$ 与 $\tau$ 都带横杠），不要读成 $\bar\partial$ 算子除以 $\partial\bar\tau$，也不要误以为这里出现了两个不同的微分算子。论证本身是标准事实：$\operatorname{Im}\tau$ 给出规范耦合、$\operatorname{Re}\tau$ 给出 $\theta$ 角，作用量对 $\bar\tau$ 求导得到 $F^+\wedge F^+$ 型（全纯）项，而它可写成 $\{Q,\cdot\}$ 的形状，于是 $Q$ 不变算符的乘积之期望值对 $\bar\tau$ 全纯。另需注意 $q=e^{2\pi i\tau}=e^{i\theta-8\pi^2/g^2}$：本丛书取 $\tau=\theta/2\pi+4\pi i/g^2$，指数上才会恰好出现 $8\pi^2/g^2$；若改用 $\tau=\theta/2\pi+i/g^2$ 的约定，这里的 $q$ 与 p.271 的瞬子数归一化会同时差一个 $4\pi$。]

理解这个求和的一个重要步骤，是确定尖点处极点的阶。人们可以有效地确定一个整数 $k_0$（一般而言可能为负），使得只要 $\mathcal M_E$ 非空就有 $k\ge k_0$。这就给出 $q=0$ 这个尖点处极点阶的一个界。把同样的估计在作一次对偶变换之后施于相应的 $k_0$，就可以确定其它尖点处极点的阶。

另一个重要步骤，是得到一个对一般的 $k$ 都成立的、简化了的计算方法。做法是注意：若 $k\gg0$，则 ${\not\partial}_A$ 的指标为负，从而至少对一般的 $A$，方程 ${\not\partial}_A(M)=0$ 除平凡解 $M=0$ 之外没有解。于是，对足够大的 $k$，$\mathcal M_E$ 上的计算就化为在更熟悉的模空间 $\mathcal M_E^{\text{instanton}}$（即 $M=0$ 的反自对偶方程之解的模空间）上的计算。此时，一个拓扑论证表明关联函数由下式给出：

\[
  \int_{\mathcal M_E^{\text{instanton}}}\chi(\mathcal V)
  \prod_i\int_{\Sigma_i^{b_i}}\mathcal O^{(b_i)},
\]

其中 $\mathcal V$ 是障碍丛（obstruction bundle），在目前情形下它等同于 ${\not\partial}_A:S^+\otimes\rho\to S^-\otimes\rho$ 的余核（cokernel），而 $\chi(\mathcal V)$ 是它的 Euler 类。（在其中一个习题中，处理了这个命题的二维类比。）

% [Tutor Note: 这里把 $\chi(\mathcal V)$ 称作「Euler 类」，而 $\chi$ 在同一页前面还表示 Euler 特征数，符号复用极易混淆：前者是障碍丛的 Euler 类（一个上同调类，次数等于丛的实秩），后者是流形的 Euler 特征数（一个数）。还要注意「障碍丛」之所以出现，是因为对每个反自对偶解，${\not\partial}_A$ 的余核给出了模空间的障碍空间；当所有瞬子光滑且不可约时，这个障碍丛恰是模空间切丛的宇称反转（参见 p.272），于是超流形上的 Berezin 积分退化为 $\mathcal M_E$ 上对 $\chi(\mathcal V)$ 的普通积分。$\mathcal V$ 的秩与 $\mathcal M_E$ 的虚维数一致，这正是「被积的是最高维形式」的来源。]

% ===== 原文 p.274 =====

有一个情形，其中对一切瞬子数，${\not\partial}_A$ 都满足消没定理，那就是：规范群为 $SU(2)$、$\rho=\mathfrak{su}(2)\otimes(\text{四元数})$、且流形为 hyper-Kähler 流形的情形。此时我们有 $\mathcal{N}=4$ 超对称，而 $\mathcal V$ 等同于瞬子模空间的切丛。特别地，$\mathcal M_E$ 的虚维数为零，所以诸 $\int_{\Sigma_i^{b_i}}\mathcal O^{(b_i)}$ 都必须取为 $1$。同时，$\chi(\mathcal V)$ 是 $\mathcal M_E$ 的 Euler 类，而在 $\mathcal M_E$ 上的积分计算的是瞬子模空间的拓扑 Euler 特征数。也就是说，我们计算的是

\[
  q^{C}\sum_k q^k\chi(\mathcal M_{E_k}),
\]

其中 $C$ 是某个适当的常数。量子场论的对偶对称性意味着：这个关于 $q$ 的级数像一个模形式那样变换（可能在尖点处有极点），其不变群是 $SL(2,\mathbb{Z})$ 的一个有限指数子群 $\Gamma$。（事实上，$\Gamma$ 就是 $2$-同余子群 $\Gamma_0(2)$。）这个预言还可以通过描述它对 $w_2(E)$ 的依赖而进一步精细化，并且已在若干细节上与数学上的计算作过成功的比较。

% [Tutor Note: 此处 $\rho=\mathfrak{su}(2)\otimes(\text{四元数})$ 意为超多重态取 $\mathfrak{su}(2)$ 的伴随表示（四元数维数为 $3=\dim G$），于是理论获得 $\mathcal{N}=4$ 超对称，且 $\mathcal{N}=4$ 的 ${\not\partial}_A$ 指标对一切 $k$ 为零——这正是 $k$ 的系数 $4h-c_2(\rho)$ 在此为零的含义，也是消没定理成立、$\mathcal V$ 恰为切丛的原因。至于虚维数整体为零，宜用 p.271 的公式在 hyper-Kähler 流形上核对而非当作恒等式：以 K3 为例，$b_1=0$、$b_2^+=3$、$\chi=24$、$\sigma=-16$，于是 $-\tfrac12\dim G(\chi+\sigma)=-3(1+b_2^+)=-12$，而 $-\tfrac{\sigma}{4}\text{qdim}(\rho)=-\tfrac{-16}{4}\cdot3=+12$，二者相消。可见「虚维数为零」依赖 hyper-Kähler 这个前提（它给出 $b_1=0$、$b_2^+=3$），并非对任意 $X$ 都成立。另注意 $\chi(\mathcal M_{E_k})$ 前的 $q^C$ 中的 $C$ 是一个与 $k$ 无关的常数（来自 $q^{a\chi+b\sigma}$ 在 $\mathcal{N}=4$ 情形的取值），不要把它误当成某个瞬子数。]

\subsubsection{Donaldson 理论}
% 原文 §19.3.2：Donaldson Theory

现在我们准备特殊化到 Donaldson 理论，也就是纯 $SU(2)$ 理论的情形；所谓「纯」，是指表示 $\rho$ 平凡，从而没有超多重态。我们把记号相对上一讲略作改动：把重整化标度 $\Lambda$ 取为 $1$，于是 $u$ 平面上的奇点位于 $u=\pm1$。（这个选择是暂时性的，将在本讲末尾略作修改。）

正如上面已指出的，基本算符是 $\mathcal O^{(0)}=u=\operatorname{Tr}\phi^2$，它来自 $\mathfrak{su}(2)$ 上的二次不变多项式；此外还有导出算符 $\mathcal O^{(i)}$，$i=1,\ldots,4$，它们要在 $X$ 中的 $i$ 维闭链上积分。我们把经典解的超流形底下那个几何流形的各个分支记作 $\mathcal M_k$，其中 $k\in\mathbb{Z}$ 是瞬子数，也就是主 $SU(2)$ 丛的第二 Chern 类的值。上面的短距离论证表明

\[
  \left\langle\prod_{i=1}^{s}\int_{\Sigma_i^{b_i}}\mathcal O^{(b_i)}\right\rangle_k
  =\int_{\mathcal M_k}\mu(\Sigma_1^{b_1})\cup\cdots\cup\mu(\Sigma_s^{b_s}),
\]

其中 $\mathcal M_k$ 是 $X$ 上第二 Chern 类为 $k$ 的瞬子模空间，只要

\[
  \mathrm{virt.\ dim}(\mathcal M_k)=\sum_{i=1}^{s}(4-b_i),
\]

或者等价地，只要

\[
  8k-\frac{3}{2}\bigl(\chi(M)-\sigma(M)\bigr)=\sum_{i=1}^{s}(4-b_i).
\]

% [Tutor Note: 此式与 p.271 的虚维数公式不一致（那里是 $\chi+\sigma$，这里成了 $\chi-\sigma$），详见 p.271 处的批注；以 $\chi+\sigma$ 为准。另外，这里的 $\chi(M)$、$\sigma(M)$ 用 $M$ 记四维流形，与全篇的 $X$ 不一致，且 $M$ 在 §19.3 中一直是超多重态的标量场，务必按 $X$ 读。]

右端的积分就是通常的 Donaldson 多项式不变量的值：

\[
  D\bigl([\Sigma_1^{b_1-1}],\ldots,[\Sigma_s^{b_2}]\bigr)
  =\int_{\mathcal M_k}\mu(\Sigma_1^{b_1})\cup\cdots\cup\mu(\Sigma_s^{b_s}).
\]

% [Tutor Note: 最后一式 $D$ 的自变量中，末项的上标印作 $b_2$，显然应为 $b_s-1$：$D$ 是关于 $s$ 个同调类的多重线性函数，若末项指数回到 $b_2$，则在 $s>2$ 时根本无法解释，且与前面 $s$ 个因子 $\mu(\Sigma_i^{b_i})$ 的对应关系断裂。另请留意 $D$ 的自变量写作 $[\Sigma_i^{b_i-1}]$、比积分中出现的 $\Sigma_i^{b_i}$ 低一维，这是 Donaldson 理论中多项式记号的约定（与 $\mu([\Sigma])$ 的次数移位相配，参见 p.271 末行 $\mu([\Sigma])\in H^{4-k}$），抄录时不要顺手把指数改成 $b_i$。]

% ===== 原文 p.275 =====

\subsection{Donaldson 理论的低能计算}
% 原文 §19.4：Low Energy Computations for Donaldson theory

为了得到有趣的结果，我们必须把这个高能计算（我们刚刚说明它归结为通常的 Donaldson 多项式不变量）与用理论低能有效极限所做的计算相比较。就度量而言，我们设想一族单参数度量 $t^2g_0$，其中 $g_0$ 是某个固定的度量。由于能动张量具有 $T=\{Q,\ldots\}$ 的形式，而我们计算的是算符的 $Q$ 上同类之关联函数，答案与度量无关。当 $t\mapsto0$ 时，我们回到上面勾勒过的短距离（亦即高能）计算，并发现结果就是 Donaldson 多项式不变量。当 $t\mapsto\infty$ 时，我们可以用理论的真空空间来计算关联函数。

这个计算背后的基本想法已在第 1 讲 §4 中解释过，例如可参看式 (1.12)。当 $t\to\infty$、度量在一切方向上都变大时，在一个只有有限多个真空态的理论中，任何关联函数的计算都将收敛为不同真空态之贡献的线性组合。在我们目前的例子中，真空态有无穷多个，由复变量 $u$ 参数化。对真空求和（在真空只有有限多个的情形）的类比，就是对 $u$ 积分。然而我们必须预料到：在这个积分中，$u=\pm1$ 处的奇点会带来 delta 函数型的贡献。理由是：在 $u=\pm1$ 处存在一些低能位形（其要害者，正如我们将看到的，是带非零磁单极子场的 Seiberg--Witten 方程之解），它们不能良好地延拓到 $u\neq\pm1$。因此，一般的结构将是 $u$ 平面上的一个积分，加上 $u=\pm1$ 处的 delta 函数型贡献。但结果表明，若 $b_2^+(X)>1$，则 $u$ 平面上的连续积分消失，而 Donaldson 不变量仅由 $u=\pm1$ 处的 delta 函数型贡献给出。其理由将在下一段简要勾勒。当 $b_2^+(X)=1$ 时，$u$ 平面确有贡献；此时 Donaldson 不变量是三项贡献之和——与前面一样是 $u=\pm1$ 处的离散贡献，再加上 $u$ 平面上某个函数的真正积分。即使 $b_2^+(X)=1$，最有意思的拓扑信息仍然在来自 $u=\pm1$ 的贡献之中，但 $u$ 平面上的积分会进入 Donaldson 不变量的公式。

除 $u=\pm1$ 之外处处，低能理论都是纯 $U(1)$ 理论。所有相互作用都是「纯粹非重整化的」（purely unrenormalizable），它们的效果当 $t\to\infty$ 时按 $t$ 的幂次消失。另一方面，若把 $u=u_0$ 附近的理论（$u_0$ 为除 $u_0=\pm1$ 之外的任意值）当作自由理论来处理，则路径积分消失，因为存在费米子零模。这些简单的观察是一个更细致分析（见本讲末尾参考文献中引用的 Moore--Witten 论文）的出发点，该分析表明：当 $b_2^+(X)>1$ 时，一般的阿贝尔真空（abelian vacua）的贡献消失。$b_2^+(X)$ 这个数之所以出现，是因为它等于前面称作 $\psi_2^+$ 的那个费米场的零模个数；这种零模越多，当 $t\to\infty$ 时就越难以避免消失。

% [Tutor Note: 「$b_2^+(X)$ 等于费米场 $\psi_2^+$ 的零模个数」一句值得展开：在扭曲后的理论中，$\psi_{(2)}$ 已是二形式场，其自对偶部分记作 $\psi_2^+$（即 $\Omega^{2,+}$ 的截面），它的零模就是调和的自对偶二形式，其个数正是 $b_2^+$；每多一个这样的零模，Grassmann 积分就多一个必须被 Yukawa 耦合「饱和」的因子，故 $b_2^+>1$ 时一般阿贝尔真空的贡献只能为零。这也解释了为何 $b_2^+=1$ 是分界：此时恰有一个零模，可以被饱和，$u$ 平面上的连续积分才不为零。另注意作者把 $u$ 平面上的连续积分与 $u=\pm1$ 处的 delta 函数型贡献并列：前者来自模空间中非紧方向上的场「跑掉」，后者来自 $u=\pm1$ 处出现的无质量超多重态，两者的物理来源不同，不可混为一谈。]

但在 $u=\pm1$ 处，情形完全不同。在这些真空中存在我们不能忽略的无质量超多重态。在 $u=\pm1$ 处，具有未破缺超对称的场位形所满足的经典方程就是 SW 方程

\[
\begin{aligned}
  F_A^+&=(M\overline M)^+,\\
  {\not\partial}_A(M)&=0.
\end{aligned}
\]

% [Tutor Note: 末式右端 $(M\overline M)^+$ 是 $M\overline M$ 的自对偶部分；这里 $M$ 是 $S^+\otimes\rho$ 的截面，$\overline M$ 属于复共轭表示，二者相乘给出 $\mathfrak{g}$ 值（或 $\mathfrak{u}(1)$ 值）的二形式，再取自对偶投影。与 p.270 的 (19.3) $F_A^+=\mu(M)$ 相比，此处把矩映射具体写成 $(M\overline M)^+$，可见 $\mu(M)$ 在这个阿贝尔真空附近就是 $M\overline M$ 的自对偶投影（相差一个由归一化决定的常数）。核对符号时注意：若把 $(M\overline M)^+$ 误写成 $(M\overline M)^-$，方程将不再有非平凡解，$u=\pm1$ 处的磁单极子贡献会整个消失。]
